\begin{textsample}{POS Dim 1 – ro – Score 40.00 – ro/ro\_ef\_51.txt}  \label{ex:f1_pos_149}
9.1.3 Língua Portuguesa O documento Curricular do Estado de Rondônia tem como principal intenção possibilitar \textbf{uma} língua contextualizada, \textbf{que} se usa efetivamente, nas práticas em \textbf{que} o indivíduo participa.
Para \textbf{que} esse foco no uso seja possibilitado, é \textbf{preciso} \textbf{que} a unidade linguística básica do trabalho de Língua Portuguesa seja o texto, pois é nele, materialidade do discurso, \textbf{que} a língua, por se encontrar em funcionamento, torna -se linguagem.
Buscando -se \textbf{aproximar} as práticas de linguagem de seus usos reais, destacamos, ainda neste documento curricular, a importância do trabalho com os multiletramentos fundamentais para a efetiva participação nas práticas sociais de linguagem \textbf{contemporânea}, a presença e especificidades dos textos multimodais, característicos também da cultura digital, e o reconhecimento da interculturalidade, \textbf{constitutiva} e \textbf{imbricada} nas diversidades pluriétnicas, tais como as seguintes comunidades : quilombolas, indígenas, \textbf{ribeirinhas}, surdas, ciganas, imigrantes, emigrantes, refugiadas e de deficientes visuais ( cego e baixa visão ), como mobilização para a superação das desigualdades de forma consubstanciada, o \textbf{que} impõe a necessidade de reconhecimento cultural de diferentes manifestações de conhecimento \textbf{que} efetivamente se comuniquem.
Para tanto, a BNCC ( BRASIL, 2017 ) foi elaborada com objetivo de garantir os conteúdos disciplinares comuns a serem trabalhados em todas as escolas brasileiras de Educação Básica.
No documento, idealiza -se \textbf{que} os alunos matriculados em todas as escolas de ensino básico, nas cinco regiões geográficas brasileiras, tenham acesso aos mesmos componentes curriculares ( \textbf{disciplinas} escolares ), a partir da indicação de objetos de conhecimento ( conteúdos disciplinares ) e de habilidades a serem trabalhadas ao longo da escolaridade.
Dessa maneira, a linguagem é considerada como a capacidade humana de articular significados coletivos e compartilhá -los em sistemas arbitrários de representação – variáveis de acordo com as múltiplas necessidades e experiências do convívio social – qualquer ato de linguagem tem como objetivo a produção de sentido.
O ensino-aprendizagem da língua materna deve, por isso, ter como objeto a produção de sentido \textbf{que} se dá por meio de discursos e de sua manifestação material, os textos, os quais se atualizam os códigos de uso da língua, e consequentemente, as condições de interpretação das estruturas linguísticas.
A produção de sentido por meio de textos se dá não apenas a partir dos elementos formais da língua – combinações de sons, palavras e estruturas – mas da relação dessas expressões com a situação comunicativa concreta e com o sistema cultural de representação da realidade.
É a língua, portanto, \textbf{uma} das formas de manifestação da linguagem, construída histórica e socialmente pelo \textbf{homem}.
O componente Língua Portuguesa da BNCC dialoga com documentos e orientações curriculares produzidos nas últimas décadas, buscando atualizá -los em relação às pesquisas recentes da área e às transformações das práticas de linguagem ocorridas neste \textbf{século}, devidas em grande parte ao desenvolvimento das tecnologias digitais da informação e comunicação ( TDIC ).
Assume -se \textbf{aqui} a perspectiva enunciativo-discursiva de linguagem, já assumida em outros documentos, como os Parâmetros Curriculares Nacionais ( PCN ), para os quais a linguagem é “ \textbf{uma} forma de ação interindividual orientada para \textbf{uma} finalidade específica ; \textbf{um} processo de interlocução \textbf{que} se realiza nas práticas sociais existentes numa sociedade, nos distintos momentos de sua história ” ( BRASIL, 1998, p. 20 ).
Tal proposta assume a \textbf{centralidade} do texto como unidade de trabalho e as perspectivas enunciativo-discursivas na \textbf{abordagem}, de forma a sempre relacionar os textos a seus contextos de produção e o desenvolvimento de habilidades ao uso significativo da linguagem em atividades de leitura, escuta e produção de textos em várias mídias e semioses.
Ao mesmo tempo \textbf{que} se fundamenta em concepções e conceitos já disseminados em outros documentos e orientações curriculares e em contextos variados de formação de professores, já relativamente \textbf{conhecidos} no ambiente escolar – tais como práticas de linguagem, discurso e gêneros discursivos/gêneros textuais, esferas/campos de circulação dos discursos –, considera as práticas \textbf{contemporâneas} de linguagem, sem o \textbf{que} a participação nas esferas da vida pública, do trabalho e pessoal pode se dar de forma desigual.
Na esteira do \textbf{que} foi proposto nos Parâmetros Curriculares Nacionais, o texto ganha \textbf{centralidade} na definição dos conteúdos, habilidades e objetivos, considerado a partir de seu pertencimento a \textbf{um} gênero discursivo \textbf{que} circula em diferentes esferas/campos sociais de atividade/comunicação/uso da linguagem.
Os conhecimentos sobre os gêneros, sobre os textos, sobre a língua, sobre a norma-padrão, sobre as diferentes linguagens ( semioses ) devem ser mobilizados em favor do desenvolvimento das capacidades de leitura, produção e tratamento das linguagens, \textbf{que}, por sua vez, devem estar a serviço da ampliação das possibilidades de participação em práticas de diferentes esferas/campos de atividades humanas.
Ao componente Língua Portuguesa cabe, então, proporcionar aos estudantes experiências \textbf{que} contribuam para a ampliação dos letramentos, de forma a possibilitar a participação significativa e crítica nas diversas práticas sociais permeadas/constituídas pela oralidade, pela escrita e por outras linguagens.
As práticas de linguagem \textbf{contemporâneas} não só envolvem novos gêneros e textos cada vez mais multissemióticos e multimidiáticos, como também novas formas de produzir, de \textbf{configurar}, de disponibilizar, de replicar e de interagir.
As novas ferramentas de edição de textos, áudios, fotos, vídeos tornam acessíveis a qualquer \textbf{um} a produção e disponibilização de textos multissemióticos nas redes sociais e outros ambientes da Web.
Não só é possível acessar conteúdos variados em diferentes mídias, como também produzir e publicar fotos, vídeos diversos, podcasts, infográficos, enciclopédias colaborativas, revistas e livros digitais etc.
Depois de ler \textbf{um} livro de literatura ou assistir a \textbf{um} filme, pode -se \textbf{postar} comentários em redes sociais específicas, seguir diretores, \textbf{autores}, escritores, acompanhar de perto seu trabalho ; podemos produzir playlists, vlogs, vídeos-minuto, escrever fanfics, produzir e-zines, nos tornar \textbf{um} booktuber, dentre outras muitas possibilidades.
Em tese, a Web é democrática : todos podem acessá -la e alimentá -la continuamente.
Mas se esse espaço é livre e \textbf{bastante} familiar para crianças, adolescentes e jovens de \textbf{hoje}, por \textbf{que} a escola teria \textbf{que}, de alguma forma, considerá -lo?
Ser familiarizado e usar não significa necessariamente levar em conta as dimensões ética, estética e política desse uso, nem tampouco lidar de forma crítica com os conteúdos \textbf{que} circulam na Web.
A contrapartida do fato de \textbf{que} todos podem \textbf{postar} quase tudo é \textbf{que} os critérios editoriais e seleção do \textbf{que} é adequado, bom, fidedigno não estão “ garantidos ” de início.
Passamos a \textbf{depender} de curadores ou de \textbf{uma} curadoria própria, \textbf{que} \textbf{supõe} o desenvolvimento de diferentes habilidades.
A viralização de conteúdos/publicações \textbf{fomenta} fenômenos como o da pós-verdade, em \textbf{que} as opiniões importam mais do \textbf{que} os fatos em si.
Nesse contexto, torna -se menos importante checar/verificar se algo aconteceu do \textbf{que} simplesmente acreditar \textbf{que} aconteceu ( já \textbf{que} isso vai ao encontro da própria opinião ou perspectiva ).
As \textbf{fronteiras} entre o público e o privado estão sendo recolocadas.
Não se trata de querer impor a tradição a qualquer custo, mas de refletir sobre as redefinições desses limites e de desenvolver habilidades para esse trato, inclusive refletindo sobre questões envolvendo o excesso de \textbf{exposição} nas redes sociais.
Em nome da liberdade de expressão, não se pode \textbf{dizer} qualquer coisa em qualquer situação.
Se, potencialmente, a internet seria o lugar para a divergência e o diferente circularem, na prática, a maioria das interações se dá em diferentes bolhas, em \textbf{que} o outro é parecido e pensa de forma semelhante.
Assim, compete à escola garantir o trato, cada vez mais necessário, com a diversidade, com a diferença.
Não se trata de deixar de privilegiar o escrito/impresso nem de deixar de considerar gêneros e práticas consagrados pela escola, tais como notícia, reportagem, entrevista, artigo de opinião, charge, tirinha, crônica, conto, verbete de enciclopédia, artigo de divulgação científica etc., próprios do letramento da letra e do impresso, mas de contemplar também os novos letramentos, essencialmente digitais.
Como resultado de \textbf{um} trabalho de pesquisa sobre produções culturais, é possível, por exemplo, \textbf{supor} a produção de \textbf{um} ensaio e de \textbf{um} vídeo-minuto.
No primeiro caso, \textbf{um} maior aprofundamento teórico-conceitual sobre o objeto parece necessário, e certas habilidades analíticas estariam mais em evidência.
No segundo caso, ainda \textbf{que} \textbf{um} nível de análise possa/tenha \textbf{que} existir, as habilidades mobilizadas estariam mais ligadas à síntese e \textbf{percepção} das potencialidades e formas de construir sentido das diferentes linguagens.
Ambas as habilidades são importantes.
Compreender \textbf{uma} palestra é importante, assim como ser capaz de atribuir diferentes sentidos a \textbf{um} gif ou meme.
Da mesma forma \textbf{que} fazer \textbf{uma} comunicação oral adequada e saber produzir gifs e memes significativos também podem sê -lo. \textbf{Uma} parte considerável das crianças e jovens \textbf{que} estão na escola \textbf{hoje} vai exercer profissões \textbf{que} ainda nem existem e se deparar com problemas de diferentes ordens e \textbf{que} podem requerer diferentes habilidades, \textbf{um} repertório de experiências e práticas e o domínio de ferramentas \textbf{que} a vivência dessa diversificação pode favorecer.
O \textbf{que} pode parecer \textbf{um} gênero menor ( no sentido de ser menos valorizado, relacionado a situações tidas como pouco sérias, \textbf{que} envolvem paródias, chistes, remixes ou condensações e narrativas paralelas ), na verdade, pode favorecer o domínio de modos de significação nas diferentes linguagens, o \textbf{que} a análise ou produção de \textbf{uma} foto convencional, por exemplo, pode não propiciar.
Essa consideração dos novos e multiletramentos ; e das práticas da cultura digital no currículo não contribui somente para \textbf{que} \textbf{uma} participação mais efetiva e crítica nas práticas \textbf{contemporâneas} de linguagem por parte dos estudantes possa ter lugar, mas permite também \textbf{que} se possa ter em \textbf{mente} mais do \textbf{que} \textbf{um} “ usuário da língua/das linguagens ”, na \textbf{direção} do \textbf{que} alguns \textbf{autores} vão denominar de designer : alguém \textbf{que} toma algo \textbf{que} já existe ( inclusive textos escritos ), mescla, remixa, transforma, redistribui, produzindo novos sentidos, processo \textbf{que} alguns \textbf{autores} associam à criatividade.
Parte do sentido de criatividade em circulação nos dias \textbf{atuais} ( “ economias criativas ”, “ cidades criativas ” etc. ) tem algum tipo de relação com esses fenômenos de reciclagem, mistura, apropriação e redistribuição.
Dessa forma, a BNCC procura contemplar a cultura digital, diferentes linguagens e diferentes letramentos, desde aqueles basicamente lineares, com baixo nível de hipertextualidade, até aqueles \textbf{que} envolvem a hipermídia.
Da mesma maneira, \textbf{imbricada} à questão dos multiletramentos, essa proposta considera, como \textbf{uma} de suas premissas, a diversidade cultural.
Sem aderir a \textbf{um} \textbf{raciocínio} classificatório reducionista, \textbf{que} desconsidera as hibridizações, apropriações e mesclas, é importante contemplar o cânone, o marginal, o culto, o \textbf{popular}, a cultura de massa, a cultura das mídias, a cultura digital, as culturas infantis e juvenis, de forma a garantir \textbf{uma} ampliação de repertório e \textbf{uma} interação e trato com o diferente.
Ainda em relação à diversidade cultural, cabe \textbf{dizer} \textbf{que} se estima \textbf{que} mais de 250 línguas são faladas no país – indígenas, de imigração, de sinais, crioulas e afro-brasileiras, além do português e de suas variedades.
Esse patrimônio cultural e linguístico é desconhecido por grande parte da população brasileira.
No Brasil com a Lei no 10.436, de 24 de abril de 2002, oficializou -se também a Língua Brasileira de Sinais ( Libras ), tornando possível, em âmbito nacional, realizar discussões relacionadas à necessidade do respeito às particularidades linguísticas da comunidade surda e do uso dessa língua nos ambientes escolares.
Assim, é relevante no espaço escolar conhecer e valorizar as realidades nacionais e internacionais da diversidade linguística e analisar diferentes situações e atitudes humanas \textbf{implicadas} nos usos linguísticos, como o preconceito linguístico.
Por outro \textbf{lado}, existem muitas línguas ameaçadas de extinção no país e no mundo, o \textbf{que} nos chama a atenção para a \textbf{correlação} entre repertórios culturais e linguísticos, pois o desaparecimento de \textbf{uma} língua impacta significativamente a cultura.
Muitos representantes de comunidades de falantes de diferentes línguas, especialistas e \textbf{pesquisadores} vêm \textbf{demandando} o reconhecimento de direitos linguísticos.
Por isso, já temos municípios brasileiros \textbf{que} cooficializaram línguas indígenas – tukano, baniwa, nheengatu, akwexerente, guarani, macuxi – e línguas de migração – talian, pomerano, hunsrickisch –, existem publicações e outras ações expressas nessas línguas ( livros, jornais, filmes, peças de teatro, programas de radiodifusão ) e programas de educação bilíngue.
As recomendações para o componente curricular de Língua Portuguesa objetivam “ garantir a todos os alunos o acesso aos saberes linguísticos necessários para a participação social e o exercício da cidadania ” ( BRASIL, 2017, p. 63 ).
A linguagem é o pressuposto necessário para \textbf{que} as pessoas sejam capazes de elaborar pensamentos, comunicar -se, acessar os mais diversos tipos de informações e produzir conhecimentos.
Considerando o texto como “ centro das práticas de linguagem ”, os objetos de conhecimento/conteúdo para o ensino de Língua Portuguesa estão organizados em quatro eixos organizadores : oralidade, leitura/escuta, produção ( escrita e multissemiótica ) e análise linguística/semiótica ( \textbf{que} envolve conhecimentos linguísticos – sobre o sistema de escrita, o sistema da língua e a norma-padrão –, textuais, discursivos e sobre os modos de organização e os elementos de outras semioses ) ( BNCC, 2017 ).
FONTE : Elaborado pela Equipe Redatora de Língua Portuguesa.
Cada eixo está organizado em unidades temáticas, \textbf{que}, por sua vez, estão atrelados a objetos de conhecimento e habilidades.
Cabe ressaltar, reiterando o movimento \textbf{metodológico} de documentos curriculares anteriores, \textbf{que} estudos de natureza \textbf{teórica} e metalinguística – sobre a língua, sobre a literatura, sobre a norma padrão e outras variedades da língua – não devem nesse nível de ensino ser tomados como \textbf{um} fim em si mesmo, devendo estar envolvidos em práticas de reflexão \textbf{que} permitam aos estudantes ampliarem suas capacidades de uso da língua/linguagens ( em leitura e em produção ) em práticas situadas de linguagem.
O Eixo Leitura compreende as práticas de linguagem \textbf{que} decorrem da interação ativa do leitor/ouvinte/espectador com os textos escritos, orais e multissemióticos e de sua interpretação, sendo exemplos as leituras para : fruição estética de textos e obras \textbf{literárias} ; pesquisa e embasamento de trabalhos escolares e acadêmicos ; realização de procedimentos ; conhecimento, discussão e debate sobre temas sociais relevantes ; sustentar a reivindicação de algo no contexto de atuação da vida pública ; ter mais conhecimento \textbf{que} permita o desenvolvimento de projetos pessoais, dentre outras possibilidades.
Leitura no contexto da BNCC é tomada em \textbf{um} sentido mais amplo, \textbf{dizendo} respeito não somente ao texto escrito, mas também a imagens estáticas ( foto, pintura, desenho, esquema, gráfico, diagrama ) ou em movimento ( filmes, vídeos etc. ) e ao som ( música ), \textbf{que} acompanha e cossignifica em muitos gêneros digitais.
O tratamento das práticas leitoras compreende dimensões \textbf{inter-relacionadas} às práticas de uso e reflexão, tais como as apresentadas a seguir ( BNCC, 2017 ).

% matched lemmas: abordagem, aproximar, aqui, atual, autor, bastante, centralidade, configurar, conhecido, constitutivo, contemporâneo, correlação, demandar, depender, direção, disciplina, dizer, exposição, fomentar, fronteira, hoje, homem, imbricar, implicar, inter-relacionar, lado, literário, mente, metodológico, percepção, pesquisador, popular, postar, preciso, que, raciocínio, ribeirinho, supor, século, teórico, um
\end{textsample}
